\documentclass[10pt,a4paper]{amsart}
\usepackage[margin=19mm]{geometry}
\usepackage{amsmath,amssymb,mathtools}
\usepackage[scr=rsfs,cal=euler]{mathalfa}
\usepackage{xfrac}
\usepackage[unicode,hidelinks,bookmarks=false]{hyperref}
\newcommand{\ie}{i.e.\ }
\newcommand{\cf}{cf.\ }
\newcommand{\resp}{respectively\ }
\newcommand{\Subsection}{\subsection}
\providecommand{\nobreakdash}{\nobreak\hbox{-}}
\setlength{\emergencystretch}{2em}
\multlinegap0pt
\usepackage{bm}
\usepackage{bbm}
\usepackage{dsfont}
\usepackage{xspace}
\usepackage{cancel}
\usepackage{mathtools}
\usepackage{amsthm}
\makeatletter
\@ifundefined{Theorem}{\newtheorem{Theorem}{Theorem}}{}
\@ifundefined{Lemma}{\newtheorem{Lemma}[Theorem]{Lemma}}{}
\@ifundefined{Remark}{\newtheorem{Remark}[Theorem]{Remark}}{}
\makeatother
\def\Aut{\mathrm{Aut}}   \def\bAut{\mathbf{Aut}}        \def\tenk{\otimes_k}     
\def\CC{{\mathcal{C}}}    \def\OP{{\mathcal{O}P}}
\def\CE{{\mathcal{E}}}    \def\OR{{\mathcal{O}R}}
\def\CL{{\mathcal{L}}}  \def\bCL{{\mathbf{L}}}
\def\Id{\mathrm{Id}}             \def\tenA{\otimes_A}
\def\N{{\mathbb N}}
\def\sym{\mathfrak{S}} %symmetric group
\title[The operad of Latin hypercubes]{The operad of Latin hypercubes}
\author{Markus Linckelmann}
\date{Source: arXiv:2402.19077v2 (CC BY 4.0)}
\begin{document}
\maketitle

\noindent\emph{Source and licence.} The operad of Latin hypercubes. Version arXiv:2402.19077v2 (CC BY 4.0), distributed under the Creative Commons Attribution 4.0 International licence, as recorded in the arXiv licence metadata for this exact version and shown on the abstract page https://arxiv.org/abs/2402.19077v2. Full rights notice: licenses/arxiv-2402.19077-CC-BY-4.0.txt.\par\smallskip

\noindent\textit{Editorial scope.} Counted unit: the Theorem whose source label is \texttt{thm2} in the article (source0.tex lines 458--465, source label \texttt{thm2}), delivered together with the article's own complete original proof of it (source0.tex lines 467--549, one \texttt{proof} environment of 83 lines). No mathematics is written, repaired or re-derived: statement and proof are the article's own text line for line. Required context retained uncounted: thm1 (lines 218--221); AutLatinRemark (lines 245--255); Lem-closed (lines 295--300); rem1 (lines 353--363). Every definition, display and label of those lines is retained unchanged. Omitted from this package and not redistributed: 1--217, 222--244, 256--294, 301--352, 364--457, 466--466, 550--971. Nothing inside a retained excerpt is omitted or altered: every retained body is delivered exactly as the article's lines are written, so no omission receipt of any kind accompanies this package. Citations: the retained excerpts carry no citation command at all, so no bibliography entry of the article is transcribed and no reference is redistributed. Reference adaptation: none was needed and none was made; every reference inside the delivered unit resolves inside it, so no unresolved reference or citation is produced. Printed slips kept literal: no printed word, symbol, hypothesis or number of the counted statement or of its proof was altered. Layout and class: the article's own class is \texttt{amsart} with packages including amssymb, amsthm, dsfont, xy; the wrapper is the lane's pinned \texttt{amsart} editorial wrapper at 10pt on A4, which restates the theorem-like environments the retained text needs and adds the article's own macro definitions that the retained text needs (\texttt{\textbackslash Aut}, \texttt{\textbackslash CC}, \texttt{\textbackslash CE}, \texttt{\textbackslash CL}, \texttt{\textbackslash Id}, \texttt{\textbackslash N}, \texttt{\textbackslash sym}), with extra wrapper packages (bm, bbm, dsfont, xspace, cancel, mathtools). The delivered numbering is the unit's own. Assets: no retained span contains a figure environment, a graphics-inclusion command, a PDF-page inclusion, a table or a TikZ picture; no image file is copied into this package, the package declares no asset dependency and no image file is redistributed with it. Licence: version arXiv:2402.19077v2 of this article is distributed under the Creative Commons Attribution 4.0 International licence, as recorded in the arXiv licence metadata for this exact version and shown on the abstract page https://arxiv.org/abs/2402.19077v2. A full rights notice is delivered at licenses/arxiv-2402.19077-CC-BY-4.0.txt. Characters: every retained excerpt is pure ASCII, so the delivered engine, XeLaTeX with its default Latin Modern text font, reports no missing character.
\begin{Theorem} \label{thm1}
Let $X$ be a non-empty set. There is a sub-operad $\CL$ of the endomorphism
operad $\CE$ of $X$ such that $\CL(d)=$ $\CL(X^d,X)$, for all  $d\in \N$.
\end{Theorem}

\begin{Remark} \label{AutLatinRemark}
Let $X$ be an object in a category $\CC$ with finite products and let $d$ be a 
positive integer.  The group $\Aut_\CC(X)\wr \sym_{d+1}$ acts on $X^{d+1}$, hence 
on the class of Latin hypercubes of  dimension $d$ over $X$, with the base group of 
$d+1$ copies of $\Aut_\CC(X)$ acting by composing $\lambda$ with a $(d+1)$-tuple 
of automorphisms of $X^{d+1}$ and $\sym_{d+1}$ acting on $X^{d+1}$ by permuting 
the $d+1$ canonical  projections  $\pi^{d+1}_i : X^{d+1}\to X$. This action induces an 
action of $\Aut_\CC(X)\wr \sym_{d+1}$ on the isomorphism classes of Latin 
hypercubes of dimension $d$ over $X$, which for $\CC$ the category of sets is
the standard notion of paratopism.
\end{Remark}

\begin{Lemma}  \label{Lem-closed}
Let $X$ be a non-empty set, and 
let $d$, $e$, $i$ be positive integers such that $1\leq i\leq d$. 
Let $f\in\CL(X^d,X)$ and $g\in\CL(X^e,X)$. 
Then $f\circ_i g\in$ $\CL(X^{d+e-1},X)$.
\end{Lemma}

\begin{Remark} \label{rem1} 
We could have very slightly simplified the proof of Lemma \ref{Lem-closed}
by observing that thanks to the symmetric group actions on the coordinates 
of  Latin hypercubes  (cf. Remark \ref{AutLatinRemark}) it would have been 
sufficient in the proof of Lemma \ref{Lem-closed}  to consider the map 
$f\circ_{d} g$ and a single $s$ in the  distinction into the  two cases  for $s$.  
That is,  it would have been sufficient in the last part of the proof of 
Lemma \ref{Lem-closed}  to consider the cases where  either $1=s\leq d-1$ or 
$s=d$.  We will make use of this observation in the proof of the more general 
Theorem \ref{thm2} below.
\end{Remark}

\begin{Theorem} \label{thm2}
Let $X$ be an object in a category $\CC$ with finite  products. 
There is a sub-operad $\CL$ of the endomorphism set
operad $\CE$ of $X$ such that $\CL(d)=$ $\CL_\CC(X^d,X)$, for all  $d\in \N$.
In particular, for any positive integers $d$, $e$, $i$ such that $1\leq i\leq d$,
and any morphisms $f\in$ $\CL_\CC(X^d,X)$ and $g\in$ $\CL_\CC(X^e,X)$ we have
$f\circ_i g \in \CL_\CC(X^{d+e-1},X)$. 
\end{Theorem}

\begin{proof}
The proof amounts to rewriting the proof of Theorem \ref{thm1}, including the statement and proof of
Lemma \ref{Lem-closed}, in such a way
that all steps remain valid for the Cartesian products in $\CC$.  In order to keep
this readable, we mention in each step what this corresponds to in the case where
$X$ is a non-empty set (and by considering coordinates, one easily translates this
to statements in $\CC$). 

\medskip
Let $f\in$ $\CL_\CC(X^d,X)$ and $g\in$ $\CL_\CC(X^e,X)$, where
$d$, $e$ are positive integers. That is, the morphisms
$$(\Id_{X^d}, f) : X^d \to X^{d+1}, $$
$$(\Id_{X^e},g) : X^e \to X^{e+1}$$
are Latin hypercubes. 
As in the proof of Theorem \ref{thm1}, the unitality and symmetric group actions
are obvious, and we only need to show, analogously to Lemma \ref{Lem-closed}, 
that  $(\Id_{X^{d+e-1}}, f\circ_i g)$ is a Latin hypercube. That is, we need to show that for
$1\leq s\leq d+e-1$ and $1\leq i\leq d$, 
the composition $\tau^{d+e}_s \circ  (\Id_{X^{d+e-1}}, f\circ_i g)$
is an automorphism of $X^{d+e-1}$. As pointed out in Remark \ref{rem1}, since
we may permute coordinates, it suffices to do this for $i=d$ and  in the two cases where
either $1=s\leq d-1$ or $s=d$.

\medskip
We consider first the case $1=s\leq d-1$, so $d\geq 2$. We need to show that the morphism
$$\tau^{d+e}_1\circ (\Id_{X^{d+e-1}}, f\circ_d g)$$
is an automorphism of $X^{d+e-1}$. If $X$ is a set, then  this automorphism is given by 
the assignment
$$(x_1,x_2,..,x_{d+e-1}) \mapsto (x_2,..,x_{d+e-1}, f(x_1,..,x_{d-1},g(x_d,..,x_{d+e-1}))).$$
First, the morphism $\tau^d_1\circ (\Id_{X^d}, f)$ is an automorphism of $X^d$ because
$(\Id_{X^d},f)$ is a Latin hypercube. We note that if $X$ is a non-empty set, then the morphism 
$\tau^d_1\circ (\Id_{X^d}, f)$ is given by the assignment 
$$(x_1,..,x_d) \mapsto (x_2,..,x_d, f(x_1,..,x_d)).$$
Compose this with the automorphism $\hat\sigma$
induced by the cyclic permutation $\sigma=$ $(1,2,..,d)$ on coordinates. 
The resulting automorphism
$$\hat\sigma\circ\tau^d_1\circ (\Id_{X^d},f)$$
is, for $X$ a set, given by the assignment
$$(x_1,..,x_d) \mapsto (f(x_1,..,x_d),x_2,..,x_d)$$
The Cartesian product  of this automorphism with $-\times \Id_{X^{e-1}}$ yields
an automorphism of $X^{d+e-1}$, which for $X$ a set corresponds to 
$$(x_1,..,x_{d+e-1}) \mapsto (f(x_1,..,x_d),x_2,..,x_{d+e-1}))$$
Again permuting cyclically the coordinates yields an automorphism of $X^{d+e-1}$ 
which we will denote by $\alpha$, which, if
$X$ is a set,  corresponds to
$$(x_1,..,x_{d+e-1}) \mapsto (x_2,..,x_{d+e-1}, f(x_1,..,x_d)).$$
Using the fact that $\tau^e_1\circ (\Id_{X^e}, g)$ is an automorphism of $X^e$,
combined with cyclically permuting coordinates, we obtain an automorphism
$\gamma$ of $X^e$ which corresponds to the assignement
$$(x_d,..,x_{d+e-1}) \mapsto (g(x_d,..,x_{d+e-1}), x_{d+1},..,x_{d+e-1}).$$
Then $\beta = \Id_{X^{d-1}}\times \eta$ is the automorphism of $X^{d+e-1}$
which corresponds to
$$(x_1,..,x_{d+e-1}) \mapsto (x_1,..,x_{d-1}, g(x_d,..,x_{d+e-1}), x_{d+1},..,x_{d+e-1}).$$
Similarly, $\gamma=\Id_{X^{d-1}}\times \eta \times \Id_X$ is an automorphism of 
$X^{d+e-1}$.
Now $\alpha\circ\beta$ is an automorphism of $X^{d+e-1}$ corresponding to
$$(x_1,..,x_{d+e-1}) \mapsto (x_2,..,x_{d-1},g(x_d,..,x_{d+e-1}),x_{d+1},..,x_{d+e-1},
f(x_1,..,x_{d-1}, g(x_d,..,x_{d+e-1}))).$$
Thus the automorphism $\gamma^{-1}\circ\alpha\circ\beta$ of $X^{d+e-1}$
coincides with the morphism $\tau^{d+e}_1\circ (\Id_{X^{d+e-1}}, f\circ_d g)$.
This proves the result in  the case $1=s\leq d-1$.

\medskip
Consider next the case $s=d$. We need to show that the morphism
$\tau^{d+e}_d\circ (\Id_{X^{d+e-1}}, f\circ_d g)$ is an automorphism of $X^{d+e-1}$.
If $X$ is a set, then this automorphism is given by the assignment
$$(x_1,..,x_{d+e-1}) \mapsto   
(x_1,..,x_{d-1},x_{d+1},..,x_{d+e-1}, f(x_1,..,x_{d-1},g(x_d,..,x_{d+e-1}))).$$
As before, by using the automorphism $\tau^e_1\circ (\Id_{X^e}, g)$,
cyclically permuting the  coordinates and then applying $\Id_{X^{d-1}}\times-$
 we obtain anautomorphism $\delta$ of $X^{d+e-1}$, which for $X$ a set
 corresponds to
 $$(x_1,..,x_{d+e-1}) \mapsto (x_1,..,x_{d-1}, g(x_d,..,x_{d+e-1}), x_{d+1},..,x_{d+e-1}).$$
 Similarly, applying $-\times \Id_{X^{e-1}}$ to the automorphism 
 $\tau^d_d\circ (\Id_{X^d}, f)$ yields an automorphism $\epsilon$ of $X^{d+e-1}$,
 which for $X$ a set corresponds to
 $$(x_1,..,x_{d+e-1}) \mapsto (x_1,..,x_{d-1}, f(x_1,..,x_d), x_{d+1},..,x_{d+e-1}).$$
 Thus $\epsilon\circ\delta$ is the automorphism of $X^{d+e-1}$ which for $X$ a set
 corresponds to 
 $$(x_1,...,x_{d+e-1}) \mapsto (x_1,..,x_{d-1},x_{d+1},..,x_{d+e-1}, (f\circ_d g)(x_1,..,x_{d+e-1})),$$
 and this is indeed the automorphism $\tau^{d+e-1}_d\circ (\Id_{X^{d+e-1}}, f\circ_d g)$.
 This proves the second case, and the result follows.
\end{proof}

\end{document}
